\IfFileExists{stacks-project.cls}{%
\documentclass{stacks-project}
}{%
\documentclass{amsart}
}

% For dealing with references we use the comment environment
\usepackage{verbatim}
\newenvironment{reference}{\comment}{\endcomment}
%\newenvironment{reference}{}{}
\newenvironment{slogan}{\comment}{\endcomment}
\newenvironment{history}{\comment}{\endcomment}

% For commutative diagrams we use Xy-pic
\usepackage[all]{xy}

% We use 2cell for 2-commutative diagrams.
\xyoption{2cell}
\UseAllTwocells

% We use multicol for the list of chapters between chapters
\usepackage{multicol}

% This is generally recommended for better output
\usepackage{lmodern}
\usepackage[T1]{fontenc}

% For cross-file-references
\usepackage{xr-hyper}

% Package for hypertext links:
\usepackage{hyperref}

% For any local file, say "hello.tex" you want to link to please
% use \externaldocument[hello-]{hello}
\externaldocument[introduction-]{introduction}
\externaldocument[conventions-]{conventions}
\externaldocument[sets-]{sets}
\externaldocument[categories-]{categories}
\externaldocument[topology-]{topology}
\externaldocument[sheaves-]{sheaves}
\externaldocument[sites-]{sites}
\externaldocument[stacks-]{stacks}
\externaldocument[fields-]{fields}
\externaldocument[algebra-]{algebra}
\externaldocument[brauer-]{brauer}
\externaldocument[homology-]{homology}
\externaldocument[derived-]{derived}
\externaldocument[simplicial-]{simplicial}
\externaldocument[more-algebra-]{more-algebra}
\externaldocument[smoothing-]{smoothing}
\externaldocument[modules-]{modules}
\externaldocument[sites-modules-]{sites-modules}
\externaldocument[injectives-]{injectives}
\externaldocument[cohomology-]{cohomology}
\externaldocument[sites-cohomology-]{sites-cohomology}
\externaldocument[dga-]{dga}
\externaldocument[dpa-]{dpa}
\externaldocument[sdga-]{sdga}
\externaldocument[hypercovering-]{hypercovering}
\externaldocument[schemes-]{schemes}
\externaldocument[constructions-]{constructions}
\externaldocument[properties-]{properties}
\externaldocument[morphisms-]{morphisms}
\externaldocument[coherent-]{coherent}
\externaldocument[divisors-]{divisors}
\externaldocument[limits-]{limits}
\externaldocument[varieties-]{varieties}
\externaldocument[topologies-]{topologies}
\externaldocument[descent-]{descent}
\externaldocument[perfect-]{perfect}
\externaldocument[more-morphisms-]{more-morphisms}
\externaldocument[flat-]{flat}
\externaldocument[groupoids-]{groupoids}
\externaldocument[more-groupoids-]{more-groupoids}
\externaldocument[etale-]{etale}
\externaldocument[chow-]{chow}
\externaldocument[intersection-]{intersection}
\externaldocument[pic-]{pic}
\externaldocument[weil-]{weil}
\externaldocument[adequate-]{adequate}
\externaldocument[dualizing-]{dualizing}
\externaldocument[duality-]{duality}
\externaldocument[discriminant-]{discriminant}
\externaldocument[derham-]{derham}
\externaldocument[local-cohomology-]{local-cohomology}
\externaldocument[algebraization-]{algebraization}
\externaldocument[curves-]{curves}
\externaldocument[resolve-]{resolve}
\externaldocument[models-]{models}
\externaldocument[functors-]{functors}
\externaldocument[equiv-]{equiv}
\externaldocument[pione-]{pione}
\externaldocument[etale-cohomology-]{etale-cohomology}
\externaldocument[proetale-]{proetale}
\externaldocument[relative-cycles-]{relative-cycles}
\externaldocument[more-etale-]{more-etale}
\externaldocument[trace-]{trace}
\externaldocument[crystalline-]{crystalline}
\externaldocument[spaces-]{spaces}
\externaldocument[spaces-properties-]{spaces-properties}
\externaldocument[spaces-morphisms-]{spaces-morphisms}
\externaldocument[decent-spaces-]{decent-spaces}
\externaldocument[spaces-cohomology-]{spaces-cohomology}
\externaldocument[spaces-limits-]{spaces-limits}
\externaldocument[spaces-divisors-]{spaces-divisors}
\externaldocument[spaces-over-fields-]{spaces-over-fields}
\externaldocument[spaces-topologies-]{spaces-topologies}
\externaldocument[spaces-descent-]{spaces-descent}
\externaldocument[spaces-perfect-]{spaces-perfect}
\externaldocument[spaces-more-morphisms-]{spaces-more-morphisms}
\externaldocument[spaces-flat-]{spaces-flat}
\externaldocument[spaces-groupoids-]{spaces-groupoids}
\externaldocument[spaces-more-groupoids-]{spaces-more-groupoids}
\externaldocument[bootstrap-]{bootstrap}
\externaldocument[spaces-pushouts-]{spaces-pushouts}
\externaldocument[spaces-chow-]{spaces-chow}
\externaldocument[groupoids-quotients-]{groupoids-quotients}
\externaldocument[spaces-more-cohomology-]{spaces-more-cohomology}
\externaldocument[spaces-simplicial-]{spaces-simplicial}
\externaldocument[spaces-duality-]{spaces-duality}
\externaldocument[formal-spaces-]{formal-spaces}
\externaldocument[restricted-]{restricted}
\externaldocument[spaces-resolve-]{spaces-resolve}
\externaldocument[formal-defos-]{formal-defos}
\externaldocument[defos-]{defos}
\externaldocument[cotangent-]{cotangent}
\externaldocument[examples-defos-]{examples-defos}
\externaldocument[algebraic-]{algebraic}
\externaldocument[examples-stacks-]{examples-stacks}
\externaldocument[stacks-sheaves-]{stacks-sheaves}
\externaldocument[criteria-]{criteria}
\externaldocument[artin-]{artin}
\externaldocument[quot-]{quot}
\externaldocument[stacks-properties-]{stacks-properties}
\externaldocument[stacks-morphisms-]{stacks-morphisms}
\externaldocument[stacks-limits-]{stacks-limits}
\externaldocument[stacks-cohomology-]{stacks-cohomology}
\externaldocument[stacks-perfect-]{stacks-perfect}
\externaldocument[stacks-introduction-]{stacks-introduction}
\externaldocument[stacks-more-morphisms-]{stacks-more-morphisms}
\externaldocument[stacks-geometry-]{stacks-geometry}
\externaldocument[moduli-]{moduli}
\externaldocument[moduli-curves-]{moduli-curves}
\externaldocument[examples-]{examples}
\externaldocument[exercises-]{exercises}
\externaldocument[guide-]{guide}
\externaldocument[desirables-]{desirables}
\externaldocument[coding-]{coding}
\externaldocument[obsolete-]{obsolete}
\externaldocument[fdl-]{fdl}
\externaldocument[index-]{index}

% Theorem environments.
%
\theoremstyle{plain}
\newtheorem{theorem}[subsection]{Theorem}
\newtheorem{proposition}[subsection]{Proposition}
\newtheorem{lemma}[subsection]{Lemma}

\theoremstyle{definition}
\newtheorem{definition}[subsection]{Definition}
\newtheorem{example}[subsection]{Example}
\newtheorem{exercise}[subsection]{Exercise}
\newtheorem{situation}[subsection]{Situation}

\theoremstyle{remark}
\newtheorem{remark}[subsection]{Remark}
\newtheorem{remarks}[subsection]{Remarks}

\numberwithin{equation}{subsection}

% Macros
%
\def\lim{\mathop{\mathrm{lim}}\nolimits}
\def\colim{\mathop{\mathrm{colim}}\nolimits}
\def\Spec{\mathop{\mathrm{Spec}}}
\def\Hom{\mathop{\mathrm{Hom}}\nolimits}
\def\Ext{\mathop{\mathrm{Ext}}\nolimits}
\def\SheafHom{\mathop{\mathcal{H}\!\mathit{om}}\nolimits}
\def\SheafExt{\mathop{\mathcal{E}\!\mathit{xt}}\nolimits}
\def\Sch{\mathit{Sch}}
\def\Mor{\mathop{\mathrm{Mor}}\nolimits}
\def\Ob{\mathop{\mathrm{Ob}}\nolimits}
\def\Sh{\mathop{\mathit{Sh}}\nolimits}
\def\NL{\mathop{N\!L}\nolimits}
\def\CH{\mathop{\mathrm{CH}}\nolimits}
\def\proetale{{pro\text{-}\acute{e}tale}}
\def\etale{{\acute{e}tale}}
\def\QCoh{\mathit{QCoh}}
\def\Ker{\mathop{\mathrm{Ker}}}
\def\Im{\mathop{\mathrm{Im}}}
\def\Coker{\mathop{\mathrm{Coker}}}
\def\Coim{\mathop{\mathrm{Coim}}}

% Boxtimes
%
\DeclareMathSymbol{\boxtimes}{\mathbin}{AMSa}{"02}

%
% Macros for moduli stacks/spaces
%
\def\QCohstack{\mathcal{QC}\!\mathit{oh}}
\def\Cohstack{\mathcal{C}\!\mathit{oh}}
\def\Spacesstack{\mathcal{S}\!\mathit{paces}}
\def\Quotfunctor{\mathrm{Quot}}
\def\Hilbfunctor{\mathrm{Hilb}}
\def\Curvesstack{\mathcal{C}\!\mathit{urves}}
\def\Polarizedstack{\mathcal{P}\!\mathit{olarized}}
\def\Complexesstack{\mathcal{C}\!\mathit{omplexes}}
% \Pic is the operator that assigns to X its picard group, usage \Pic(X)
% \Picardstack_{X/B} denotes the Picard stack of X over B
% \Picardfunctor_{X/B} denotes the Picard functor of X over B
\def\Pic{\mathop{\mathrm{Pic}}\nolimits}
\def\Picardstack{\mathcal{P}\!\mathit{ic}}
\def\Picardfunctor{\mathrm{Pic}}
\def\Deformationcategory{\mathcal{D}\!\mathit{ef}}
\usepackage{cleveref}
\crefname{originallemma}{original lemma}{original lemmas}
\Crefname{originallemma}{Original lemma}{Original lemmas}
\makeatletter
\newlabel{properties-lemma-colimit-kappa}{{077N}{source}{}{sourcepointer-1}{}}
\newlabel{properties-lemma-colimit-kappa@cref}{{[originallemma][0][]077N}{[1][0][]source}{}{}{}}
\makeatother
\crefname{originallemma}{original lemma}{original lemmas}
\Crefname{originallemma}{Original lemma}{Original lemmas}
\makeatletter
\newlabel{schemes-lemma-widetilde-pullback}{{01I9}{source}{}{sourcepointer-2}{}}
\newlabel{schemes-lemma-widetilde-pullback@cref}{{[originallemma][0][]01I9}{[1][0][]source}{}{}{}}
\makeatother
\title{Quasi-coherent modules on a groupoid scheme with flat s and t are directed colimits of their kappa-generated submodules}
\author{The Stacks Project Authors}
\date{stacks-project:groupoids.tex:local-snapshot-20261012; GNU FDL 1.2 or later}
\begin{document}
\maketitle
\setcounter{subsection}{1}
\noindent\textit{Editorial note.} All delivered mathematical text is the author's own, copied from the pinned Stacks Project chapter file groupoids.tex; no statement or proof step was rewritten, reordered or removed. The counted span is the author environment and its own proof environment. References inside the span are handed to the assembler as navigation stubs that carry the canonical tag of the cited tag only; whether each of them is load-bearing is recorded on the stub itself.\par\medskip
\noindent\textit{Allowed standard prerequisites (editorial).}\begin{itemize}
\item Lemma properties-lemma-colimit-kappa (canonical tag 077N): the scheme-level colimit statement whose proof is repeated here. cited at proof lines 3386-3388. Navigation anchor for the strategy; the counted span repeats the construction in full.
\item Lemma schemes-lemma-widetilde-pullback (canonical tag 01I9): pullback of the tilde of a module is the tilde of the base change. cited at proof lines 3413-3423. Justifies the alpha-isomorphisms on the affine pieces W\_ijk used to define the closure operators.
\end{itemize}\medskip
\ifdefined\hypertarget\hypertarget{sourcepointer-1}{}\fi
\noindent\textit{Source reference.} Original lemma 077N: sources/stacks-chow/tags.txt, source lines 9208--9208. Cited at proof lines 3386-3388 as the source of the repeated argument; the counted proof repeats the construction in full for the groupoid setting, so this is the navigation anchor for the strategy.\par\medskip
\ifdefined\hypertarget\hypertarget{sourcepointer-2}{}\fi
\noindent\textit{Source reference.} Original lemma 01I9: sources/stacks-chow/tags.txt, source lines 1970--1970. Cited at proof lines 3413-3423 to justify alpha|\_\{W\_ijk\} as the displayed isomorphism on the affine pieces; load-bearing for the descent datum of the submodule.\par\medskip
\begin{lemma}
\label{lemma-colimit-kappa}
Let $(U, R, s, t, c)$ be a groupoid scheme over $S$.
Assume that $s, t$ are flat. There exists a
cardinal $\kappa$ such that every quasi-coherent module
$(\mathcal{F}, \alpha)$ on $(U, R, s, t, c)$
is the directed colimit of its $\kappa$-generated
quasi-coherent submodules.
\end{lemma}

\begin{proof}
In the statement of the lemma and in this proof
a {\it submodule} of a quasi-coherent module $(\mathcal{F}, \alpha)$
is a quasi-coherent submodule $\mathcal{G} \subset \mathcal{F}$
such that $\alpha(t^*\mathcal{G}) = s^*\mathcal{G}$ as subsheaves of
$s^*\mathcal{F}$. This makes sense because since $s, t$ are flat the
pullbacks $s^*$ and $t^*$ are exact, i.e., preserve subsheaves.
The proof will be a repeat of the proof of
Properties, Lemma \ref{properties-lemma-colimit-kappa}.
We urge the reader to read that proof first.

\medskip\noindent
Choose an affine open covering $U = \bigcup_{i \in I} U_i$.
For each pair $i, j$ choose affine open coverings
$$
U_i \cap U_j = \bigcup\nolimits_{k \in I_{ij}} U_{ijk}
\quad\text{and}\quad
s^{-1}(U_i) \cap t^{-1}(U_j) = \bigcup\nolimits_{k \in J_{ij}} W_{ijk}.
$$
Write $U_i = \Spec(A_i)$, $U_{ijk} = \Spec(A_{ijk})$,
$W_{ijk} = \Spec(B_{ijk})$.
Let $\kappa$ be any infinite cardinal $\geq$ than the cardinality
of any of the sets $I$, $I_{ij}$, $J_{ij}$.

\medskip\noindent
Let $(\mathcal{F}, \alpha)$ be a quasi-coherent module on $(U, R, s, t, c)$.
Set $M_i = \mathcal{F}(U_i)$, $M_{ijk} = \mathcal{F}(U_{ijk})$.
Note that
$$
M_i \otimes_{A_i} A_{ijk} = M_{ijk} = M_j \otimes_{A_j} A_{ijk}
$$
and that $\alpha$ gives isomorphisms
$$
\alpha|_{W_{ijk}} :
M_i \otimes_{A_i, t} B_{ijk}
\longrightarrow
M_j \otimes_{A_j, s} B_{ijk}
$$
see
Schemes, Lemma \ref{schemes-lemma-widetilde-pullback}.
Using the axiom of choice we choose a map
$$
(i, j, k, m) \mapsto S(i, j, k, m)
$$
which associates to every $i, j \in I$, $k \in I_{ij}$ or $k \in J_{ij}$
and $m \in M_i$ a finite subset $S(i, j, k, m) \subset M_j$
such that we have
$$
m \otimes 1 = \sum\nolimits_{m' \in S(i, j, k, m)} m' \otimes a_{m'}
\quad\text{or}\quad
\alpha(m \otimes 1) = \sum\nolimits_{m' \in S(i, j, k, m)} m' \otimes b_{m'}
$$
in $M_{ijk}$ for some $a_{m'} \in A_{ijk}$ or $b_{m'} \in B_{ijk}$.
Moreover, let's agree that $S(i, i, k, m) = \{m\}$ for all
$i, j = i, k, m$ when $k \in I_{ij}$. Fix such a collection $S(i, j, k, m)$

\medskip\noindent
Given a family $\mathcal{S} = (S_i)_{i \in I}$ of subsets
$S_i \subset M_i$ of cardinality at most $\kappa$ we set
$\mathcal{S}' = (S'_i)$ where
$$
S'_j = \bigcup\nolimits_{(i, j, k, m)\text{ such that }m \in S_i}
S(i, j, k, m)
$$
Note that $S_i \subset S'_i$. Note that $S'_i$ has cardinality at most
$\kappa$ because it is a union over a set of cardinality at most $\kappa$
of finite sets. Set $\mathcal{S}^{(0)} = \mathcal{S}$,
$\mathcal{S}^{(1)} = \mathcal{S}'$ and by induction
$\mathcal{S}^{(n + 1)} = (\mathcal{S}^{(n)})'$. Then set
$\mathcal{S}^{(\infty)} = \bigcup_{n \geq 0} \mathcal{S}^{(n)}$.
Writing $\mathcal{S}^{(\infty)} = (S^{(\infty)}_i)$ we see that
for any element $m \in S^{(\infty)}_i$ the image of $m$ in
$M_{ijk}$ can be written as a finite sum $\sum m' \otimes a_{m'}$
with $m' \in S_j^{(\infty)}$. In this way we see that setting
$$
N_i = A_i\text{-submodule of }M_i\text{ generated by }S^{(\infty)}_i
$$
we have
$$
N_i \otimes_{A_i} A_{ijk} = N_j \otimes_{A_j} A_{ijk}
\quad\text{and}\quad
\alpha(N_i \otimes_{A_i, t} B_{ijk}) = N_j \otimes_{A_j, s} B_{ijk}
$$
as submodules of $M_{ijk}$ or $M_j \otimes_{A_j, s} B_{ijk}$.
Thus there exists a quasi-coherent submodule
$\mathcal{G} \subset \mathcal{F}$ with $\mathcal{G}(U_i) = N_i$
such that $\alpha(t^*\mathcal{G}) = s^*\mathcal{G}$ as submodules
of $s^*\mathcal{F}$. In other words,
$(\mathcal{G}, \alpha|_{t^*\mathcal{G}})$ is a submodule of
$(\mathcal{F}, \alpha)$.
Moreover, by construction $\mathcal{G}$ is $\kappa$-generated.

\medskip\noindent
Let $\{(\mathcal{G}_t, \alpha_t)\}_{t \in T}$ be the set of
$\kappa$-generated quasi-coherent submodules of $(\mathcal{F}, \alpha)$.
If $t, t' \in T$ then $\mathcal{G}_t + \mathcal{G}_{t'}$ is also a
$\kappa$-generated quasi-coherent submodule as it is the image of the map
$\mathcal{G}_t \oplus \mathcal{G}_{t'} \to \mathcal{F}$.
Hence the system (ordered by inclusion) is directed.
The arguments above show that every section of $\mathcal{F}$ over $U_i$
is in one of the $\mathcal{G}_t$ (because we can start with $\mathcal{S}$
such that the given section is an element of $S_i$). Hence
$\colim_t \mathcal{G}_t \to \mathcal{F}$ is both injective and surjective
as desired.
\end{proof}

\end{document}
